% Generated draft: requires mathematical and source review.
\documentclass[11pt]{article}
\usepackage[margin=25mm]{geometry}
\usepackage{fontspec}
\setmainfont{Latin Modern Roman}
\usepackage{amsmath,amssymb,mathtools,unicode-math}
\setmathfont{Latin Modern Math}
\usepackage{enumitem,needspace,xurl,hyperref}
\setlength{\parindent}{0pt}
\setlength{\parskip}{5pt}
\setlength{\emergencystretch}{3em}
\newcommand{\N}{\mathbb{N}}
\newcommand{\Z}{\mathbb{Z}}
\newcommand{\Q}{\mathbb{Q}}
\newcommand{\R}{\mathbb{R}}
\newcommand{\C}{\mathbb{C}}
\newcommand{\F}{\mathbb{F}}
\newcommand{\A}{\mathbb{A}}
\newcommand{\PP}{\mathbb{P}}
\newcommand{\E}{\mathbb{E}}
\newcommand{\eps}{\varepsilon}
\newcommand{\dd}{\,\mathrm{d}}
\begin{document}
\section*{UnsolvedMath ID: 30005182}
\textbf{Source problem code:} catalogue entry
\subsection{Definitions and mathematical statement}
A pair of smooth four-manifolds \(X_0,X_1\) is an \emph{exotic pair} if their underlying topological manifolds are homeomorphic but \(X_0\) and \(X_1\) are not diffeomorphic. The supplied passage uses this absolute notion of exoticness; it does not specify equivalences relative to a chosen boundary identification.

The reconstructed request is to construct exotic smooth four-manifolds with a fixed boundary three-manifold and fixed topological classification data. A conditional template for such a request is the following. Given a specified class \(\mathcal C\) of smooth four-manifolds with boundary, a specified three-manifold \(M\), and specified invariant data \(I\) with a prescribed value \(v\), seek \(X_0,X_1\in\mathcal C\) satisfying
\[
\partial X_0\cong_{\mathrm{homeo}}M\cong_{\mathrm{homeo}}\partial X_1,
\qquad I(X_0)=I(X_1)=v,
\]
\[
X_0\cong_{\mathrm{homeo}}X_1,
\qquad X_0\not\cong_{\mathrm{diff}}X_1.
\]
Here \(I\) is a placeholder for the intended collection of invariants. This is only a template: the source does not specify \(\mathcal C\), \(M\), \(I\), \(v\), or whether these parameters are to be chosen existentially or prescribed universally. No assertion of existence for every choice of these parameters is justified. Equality of genuine homeomorphism invariants already follows from the displayed homeomorphism; requiring prescribed values would impose additional restrictions.

The supplied contextual passage discusses four-manifolds with \(\pi_1(X)\cong\{1\}\) or \(\pi_1(X)\cong\Z\), both with and without boundary. It describes the equivariant intersection form using the integral homology of the universal cover \(\widetilde X\), and says that in the simply connected case this reduces to the intersection form on \(H_2(X;\Z)\). It also mentions the boundary homeomorphism type and two additional, unnamed invariants relating the algebraic topology of \(X\) to that of \(\partial X\). Their definitions and the hypotheses needed to specify the intersection form precisely are absent. In particular, the excerpt does not establish whether the intended construction is restricted to either of the two displayed fundamental groups.

Consequently, a uniquely quantified mathematical construction problem cannot be recovered from this record. Compactness, connectedness, orientability, boundary regularity conventions, the strength of the fixed boundary identification, and the required number of distinct smooth structures remain unspecified.
\subsection{Short English statement}
The apparent aim is to construct homeomorphic but nondiffeomorphic smooth four-manifolds while prescribing their boundary and certain topological invariants. The supplied excerpt gives context rather than a complete construction problem: it does not identify the boundary, the invariant values, or the necessary quantifiers and hypotheses. The displayed formulation is therefore a conditional template requiring human mathematical review.
\subsection{Sources}
\begin{itemize}
\item Supplied local UnsolvedMath registry record, permanent ID \(30005182\), title ``Controlled Exotic Smooth Four-Manifolds with Boundary,'' exact source problem code \texttt{catalogue entry}. No proposer is identified; the proposed year is recorded as \(2022\).
\item Supplied source citation: \emph{Topologie (2023)}, Oberwolfach Reports, DOI \texttt{10.4171/owr/2022/34}, pp.\ 1927--1986. Source URL and supplied literature reference labelled ``original Oberwolfach report'': \url{https://doi.org/10.4171/owr/2022/34}. This is the supplied authoritative reference; its contents have not been independently consulted here.
\item Literature assessment dated exactly \texttt{2026-08-22}: the supplied curation retains the record as open because the extracted text is a contextual fragment and does not identify a unique classification claim that can be checked. This is supplied dated curation, not an independently verified assessment of the literature or current status.
\item The registry labels the statement \texttt{corrected\_verified} and says its reconstruction was checked against the official report. These are supplied verification claims, not verification performed for this draft.
\item The excerpt cites report-local references [12], [11], [18], [4], and [6] for classification results; [13] for some finite cyclic groups; [14] for solvable Baumslag--Solitar groups; [5] for further context; and [17] for some smooth exotica results with arbitrary fundamental group. The corresponding bibliography entries and theorem hypotheses were not supplied, so these references cannot be identified more precisely.
\end{itemize}
\end{document}
